\section{Compression and release with inherited memory}
\label{sec:contact}
The first touchdown may have positive, zero or negative $p$; subsequent
touchdowns have $p<0$.  Their contact passages have the same two time
scales.  Compression reduces the incoming $O(\delta)$ slip to
$O(\eps)$ on the time scale $\eps/\sqrt\delta$; the normal motion
then releases on the time scale $\eps$.  The friction tail contributes
different errors to velocity and displacement, so we estimate its
first two time moments.  A projection of the normal velocity removes
the stiff slip derivative during release and carries the displacement
and memory inherited from compression.  The comparison retains the
finite incoming scale, including the case $p_{\rm in}=0$.

\subsection{Incoming data and the finite-scale reference}

Write $u_r=u_0-r$, where the belt offset $r$ is constant during the
passage.  The determinant of $(p,u_0)$ with respect to $(\phi,\omega)$
at $G$ is $a\cos\phi^*=35/136$.  Thus $(p,u_0,y,v)$ is another regular
mechanical chart, on the same local branch as the chart in
Section~\ref{sec:model}.  In particular, $b$ is a dependent function of
these four coordinates.

Fix $C_0,C_{\rm in},Z_*,\gamma_0>0$ and a compact set
\begin{equation}
 \mathcal K\Subset\{(\pi,Q,S):Q>0,\ S>0,\ \pi S<2hQ\}.
 \label{con:input-compact}
\end{equation}
All constants below are uniform on these sets and on the mechanical
and profile classes of Section~\ref{sec:model}.  We consider actual
physical inputs at time $t_{\rm in}$ satisfying
\begin{equation}
 \begin{gathered}
 p_{\rm in}=\pi\delta,\quad v_{\rm in}=-\delta^2Q,\quad
 u_{r,\rm in}=\delta S,\quad
 y_{\rm in}=z_{\rm in}=\eps\delta^2\zeta,\\
 (\pi,Q,S)\in\mathcal K,\qquad 0<\eps\le C_0\delta^2,\\
 |\zeta|\le Z_*,\quad Q-\zeta\ge\gamma_0,\quad
 |X_{\rm in}-G|\le C_{\rm in}\delta.
 \end{gathered}
 \label{con:input}
\end{equation}
The inverse uses $u_0=u_r+r$; $b_{\rm in}$ is evaluated at this same
state.  The mechanical localization in \eqref{con:input} implies
$b_{\rm in}=O(\delta)$.  It allows $r=O(\delta)$, as needed at the
first touchdown.  The smaller offset range used for subsequent returns
will be imposed only in that application.

Put
\begin{equation}
 \sigma=\sqrt\delta,\qquad e=\eps\delta^{-3/2},\qquad
 \tau=e\sigma=\eps/\delta,\qquad
 h_{\rm in}=\mu A(\phi_{\rm in})-B(\phi_{\rm in}).
 \label{con:scales}
\end{equation}
Then $e\le C_0\sigma$, with no positive lower bound on $e$.
Use compression time $s=\sigma(t-t_{\rm in})/\eps$ and variables
\begin{equation}
 (y,z)=\eps\delta^{3/2}(Y,Z),\qquad v=\delta^2V_c,
 \qquad u_r=\delta U,\qquad L=Z-Y.
 \label{con:compression-chart}
\end{equation}
The saturated reference is the signed linear system
\begin{equation}
 \begin{gathered}
 \bar Y'=\bar V_c,\qquad \sigma\bar Z'=\bar Y-2\bar Z,
 \qquad \bar L=\bar Z-\bar Y,\\
 \bar V_c'=\pi\bar L,\qquad \bar U'=-h_{\rm in}\bar L,\\
 (\bar Y,\bar Z,\bar V_c,\bar U)(0)
       =(\sigma\zeta,\sigma\zeta,-Q,S).
 \end{gathered}
 \label{con:reference}
\end{equation}
We shall prove that its first arrest $T_\sigma$, defined by $\bar U=0$,
occurs in positive force.  From that event the inherited normal
reference is
\begin{equation}
 N_0'=A_DN_0,\quad
 A_D=\begin{pmatrix}0&1&0\\-1/2&0&1/2\\1&0&-2\end{pmatrix},
 \qquad N_0(0)=(\bar Y,\sigma\bar V_c,\bar Z)(T_\sigma).
 \label{con:inherited-release}
\end{equation}
Here $N_0=(Y_0,V_0,Z_0)$ and its elapsed time is denoted by $\rho$.
Its first force zero is $R_0$; its initial memory is the value
$\bar Z(T_\sigma)$ inherited from compression.

The zero-scale compression solves $y''=-\pi y/2$, $u'=hy/2$,
with $(y,y',u)(0)=(0,-Q,S)$.  Its first arrest time and amplitude are
\begin{equation}
 \begin{gathered}
 \xi_s=\begin{cases}
 \sqrt{2/\pi}\arccos(1-\pi S/(hQ)),&\pi>0,\\
 2\sqrt{S/(hQ)},&\pi=0,\\
 \sqrt{2/(-\pi)}\operatorname{arcosh}(1+(-\pi)S/(hQ)),&\pi<0,
 \end{cases}
 \qquad
 A_s^2=\frac{QS}{h}-\frac{\pi S^2}{2h^2}>0,\\
 (y,y',u)(\xi_s)=(-2A_s,-Q+\pi S/h,0).
 \end{gathered}
 \label{con:amplitude}
\end{equation}
The formulas extend smoothly through $\pi=0$.  For $\pi>0$ the strict
inequality in \eqref{con:input-compact} says exactly that arrest
precedes the first return of $y$ to zero.

\begin{lemma}[Normal opening]\label{con:normal-opening}
Let $n_\#'=A_Dn_\#$ and $n_\#(0)=(-2,0,-1)$.  Its force has a finite
first zero $\rho_\#$, with
\begin{equation}
 Y_\#(\rho_\#)=Z_\#(\rho_\#)=-\gamma_\#<0,\qquad
 V_\#(\rho_\#)=\nu_\#>\gamma_\#>0.
 \label{con:canonical-release}
\end{equation}
\end{lemma}
\begin{proof}
While $L_\#=Z_\#-Y_\#>0$, $V_\#'=L_\#/2>0$.  If there were no
force zero, $V_\#$ would eventually have a positive lower bound and
$Y_\#$ would reach zero.  Before that point,
$Z_\#(\rho)=-e^{-2\rho}+\int_0^\rho e^{-2(\rho-t)}Y_\#(t)\,dt<0$,
which contradicts positive force at $Y_\#=0$.  Hence the first force
zero is finite, occurs with $Y_\#=Z_\#<0$, and has $V_\#>0$.
The identities
$L_\#'=-Y_\#-2L_\#-V_\#$ and
$L_\#''=-V_\#-2L_\#'-L_\#/2$ show that its derivative there is
strictly negative: if it vanished, the second identity would make
the force negative just before its first zero.  At the zero,
$-L_\#'=V_\#+Z_\#=\nu_\#-\gamma_\#$.
\end{proof}

\begin{theorem}[Contact from a complete incoming state]
\label{con:contact-release}
For all sufficiently small $\delta$, every input \eqref{con:input}
has the following contact passage. The constant $M_{\rm sl}$ is fixed
independently of $\eps,\delta$.

\smallskip\noindent
\textup{(i) Compression to $w=M_{\rm sl}$.}
Positive contact persists to the first descending section
$w=u_r/\eps=M_{\rm sl}$. At this section,
\begin{equation}
 \begin{gathered}
 L_{\rm sl}\ge\ell_a>0,\quad U'(s_{\rm sl})\le-g_a<0,\qquad
 |s_{\rm sl}-T_\sigma|\le C(e+\tau),\\
 |(Y,\sigma V_c,Z)(s_{\rm sl})-N_0(0)|\le Ce.
 \end{gathered}
 \label{con:M-admission}
\end{equation}
\smallskip\noindent
\textup{(ii) Release and opening.}
The passage then reaches its first release.
There is a common bracket $0<r_-<R_0<r_+<R_*$ for the inherited first
release.  In elapsed release time $\rho=(t-t_{\rm sl})/\eps$, set
$v=\delta^{3/2}\widehat V$.  On every positive-contact prefix through
the actual first force zero $R_\eps$,
\begin{equation}
 \begin{gathered}
 |(Y,\widehat V,Z)-N_0|+|L-L_0|+|L'-L_0'|\le Ce,
 \qquad L_0=Z_0-Y_0,\\
 |R_\eps-R_0|+
 |(Y,\widehat V,Z)(R_\eps)-N_0(R_0)|\le Ce,
 \qquad R_\eps\in(r_-,r_+).
 \end{gathered}
 \label{con:release-comparison}
\end{equation}
At release, for common positive constants,
\begin{equation}
 y_L=z_L<0,\quad v_L>0,\quad
 w_b\le u_{r,L}/\eps\le W_*,\qquad
 v_L+z_L/\eps\ge \kappa_{\rm op}\delta^{3/2}>0.
 \label{con:opening}
\end{equation}
\reservestatementspace
\smallskip\noindent
\textup{(iii) Physical state, time and leading increments.}
Reconstruct a reference physical release in the same inverse chart by
\begin{equation}
 p_{L,0}=\pi\delta,\quad u_{r,L,0}=0,\quad u_{0,L,0}=r,
 \quad y_{L,0}=z_{L,0}=0,\quad
 v_{L,0}=\delta^{3/2}V_0(R_0).
 \label{con:physical-reconstruction}
\end{equation}
With $\Delta$ denoting actual minus reconstructed values,
\begin{equation}
 |\Delta p_L|+|\Delta b_L|\le Ce\delta,\qquad
 |\Delta u_{r,L}|+|\Delta v_L|\le Ce\delta^{3/2},\qquad
 |y_L|+|z_L|\le Ce\delta^3.
 \label{con:physical-errors}
\end{equation}
For $t_{L,0}=t_{\rm in}+(\eps/\sigma)T_\sigma+\eps R_0$,
\begin{equation}
 |t_{L,\eps}-t_{L,0}|
 \le C\{(\eps/\sigma)(e+\tau)+\eps e\}\le C(C_0)\eps.
 \label{con:release-clock}
\end{equation}
The entire physical contact prefix satisfies
\begin{equation}
 \begin{gathered}
 t_L-t_{\rm in}\le C\eps/\sqrt\delta,\quad
 \sup|\cX-G|\le C\delta,\quad
 \sup|p-\pi\delta|\le Ce\delta,\quad \dot p<-a/2,\\
 \sup(|y|+|z|)\le C\eps\delta^{3/2},\quad
 \sup|v|\le C\delta^{3/2},\qquad
 0\le\int_{t_{\rm in}}^t\lambda_N\,dt\le C\delta.
 \end{gathered}
 \label{con:prefix}
\end{equation}
Finally, using the actual normalized inputs in \eqref{con:amplitude},
\begin{equation}
 \begin{aligned}
 v_L&=\delta^{3/2}\{A_s\nu_\#+O(\sigma)\},&
 y_L=z_L&=\eps\delta^{3/2}\{-A_s\gamma_\#+O(\sigma)\},\\
 q_L-q_{\rm in}&=(\eps/\sigma)
       \{\xi_s\dot q_{\rm in}+O(\sigma)\},&
 p_L-p_{\rm in}&=(\eps/\sigma)\{-a\xi_s+O(\sigma)\},\\
 \dot q_L-\dot q_{\rm in}&=\delta(S/h)(-5/4,1)^T+O(\delta^{3/2}),&
 b_L-b_{\rm in}&=a\beta\delta S/h+O(\delta^{3/2}).
 \end{aligned}
 \label{con:leading-data}
\end{equation}
These estimates compare values of the actual compliant event map.
\end{theorem}

\subsection{Proof of the contact estimates}

\begin{proof}[Proof of Theorem~\ref{con:contact-release}]
We first obtain positive-force margins for the reference, then compare
the physical compression up to a fixed small slip and along its
descending tail.  The tail force enters velocity directly and position
after one time integration; both of its time moments are needed.
We then project out the stiff slip derivative during release, retaining
the inherited displacement and memory, and reconstruct the physical
state and its leading increments.

\paragraph{Step 1: Positive reference compression.}
The reference filter gives the exact formulas
\begin{equation}
 \begin{aligned}
 R_f:=\bar Z-\bar Y/2
 &=\tfrac12\sigma\zeta e^{-2s/\sigma}
   -\tfrac12\int_0^s e^{-2(s-t)/\sigma}\bar V_c(t)\,dt,\\
 \bar L(s)&=-\tfrac12\sigma\zeta(1-e^{-2s/\sigma})
 -\tfrac12\int_0^s(1+e^{-2(s-t)/\sigma})\bar V_c(t)\,dt.
 \end{aligned}
 \label{con:filter}
\end{equation}
The kernel mass is at most $\sigma/2$.  Its integral form and
Gronwall's inequality bound the reference on every fixed finite
interval, uniformly in $\sigma$, and give $R_f=O(\sigma)$.
Substituting $\bar L=-\bar Y/2+R_f$ gives an $O(\sigma)$ comparison
with the zero-scale system in \eqref{con:amplitude}.  Compactness of
$\mathcal K$ gives $0<T_-\le\xi_s\le T_+$, $A_s\ge A_->0$, and a
common positive-force neighborhood of arrest.

Positivity close to the initial force zero needs a separate estimate.
Equation~\eqref{con:filter} first gives $|\bar L(s)|\le Cs$, hence
$|\bar V_c(s)+Q|\le Cs^2$.  Choose a common small $s_0$ such that
$\bar V_c\le-Q+\eta$ there, with
$\eta\le\min(Q_-/2,\gamma_0/2)$.  The second formula then yields
\begin{equation}
 \bar L(s)\ge
 \min\{(Q-\eta)/2,Q-\eta-\zeta\}s
 \ge\min\{Q_-/4,\gamma_0/2\}s\quad(0<s\le s_0).
 \label{con:initial-positive}
\end{equation}
Indeed, a negative coefficient of $\sigma(1-e^{-2s/\sigma})$ is
bounded using $1-e^{-2s/\sigma}\le2s/\sigma$; a positive one is
discarded.  Beyond $s_0$ the reduced positive force and the
$O(\sigma)$ comparison exclude a force zero up to a fixed extension
past arrest.  Thus $\bar L\ge\kappa_Ls$ on the required interval, for a common
$\kappa_L>0$, and
$\bar U'=-h_{\rm in}\bar L$ gives a unique first arrest with
$T_\sigma=\xi_s+O(\sigma)$.

The exact invariant gives
$\bar V_c(T_\sigma)=-Q+\pi S/h_{\rm in}$.  An amplitude calculation
which does not divide by $\pi$ is useful here.  Put
$I(s)=\int_0^s\bar L\,dt$, so that
$\bar V_c=-Q+\pi I$, $I'=-\bar Y/2+R_f$ and
$I(T_\sigma)=S/h_{\rm in}$.  Integrating
$2\bar Y\bar V_c=-4\bar V_c I'+4\bar V_cR_f$ gives
\begin{equation}
 \bar Y(T_\sigma)^2=\sigma^2\zeta^2+
 \frac{4QS}{h_{\rm in}}-\frac{2\pi S^2}{h_{\rm in}^2}
 +4\int_0^{T_\sigma}\bar V_cR_f\,ds.
 \label{con:amplitude-identity}
\end{equation}
The force margin selects its negative square root.  Since
$h_{\rm in}=h+O(\delta)$,
\begin{equation}
 N_0(0)=A_s(-2,0,-1)+O(\sigma).
 \label{con:inherited-amplitude}
\end{equation}
Choose $0<u_a<S_-/4$ small enough that the entire reference segment
$-u_a/2\le\bar U\le u_a$ lies in the above arrest neighborhood and
has $\bar L\ge4\ell_a>0$, $-\bar U'\ge m_a>0$.  This choice uses
the compact strict margin in \eqref{con:input-compact}.
The fixed profile conditions give, on the mechanical chart,
\begin{equation}
 H_\theta(w,X):=\mu A\Phi_\theta(w)-B
 \ge\kappa_H:=\mu A_{\min}\varkappa_{\rm sl}>0
 \quad(w\ge M_{\rm sl}).
 \label{con:tail-margin}
\end{equation}
Here $A_{\min}>0$ is the common lower bound of $A$ on the chart;
monotonicity and (F3) give
$\Phi_\theta(w)-B/(\mu A)\ge\varkappa_{\rm sl}$.

\paragraph{Step 2: Compression to the small-slip section.}
On positive contact the physical equations give exactly, with
$F_\theta=1-\Phi_\theta$,
\begin{equation}
 \begin{aligned}
 Y'&=V_c,&\sigma Z'&=Y-2Z,\\
 V_c'&=(p/\delta)L+(e/\delta)b
       +(\mu B/\delta)F_\theta(U/\tau)L,\\
 U'&=ea_T-h(\phi)L+\mu AF_\theta(U/\tau)L,
 \qquad h(\phi)=\mu A-B.
 \end{aligned}
 \label{con:actual-compression}
\end{equation}
Here $\lambda_N=e^{-1}L$, $dt/ds=e\delta$ and
$u_r/\eps=U/\tau$.  The global memory row is instead
$\sigma Z'=-Z-[Z-Y]_+$; we use the displayed contact row only up to
its first force zero.  At the initial point $L'(0)=Q-\zeta>0$.

For mechanical increments set
$\chi=(q-q_{\rm in})/(e\delta)$ and
$\mathcal V=(\dot q-\dot q_{\rm in})/\delta$.  The exact rows are
\begin{equation}
 \chi'=\dot q_{\rm in}+\delta\mathcal V,\qquad
 \mathcal V'=eM^{-1}f+
 M^{-1}(Q_N-\mu Q_T\Phi_\theta(U/\tau))L.
 \label{con:mechanical-increments}
\end{equation}
On a bounded stopped interval with bounded $\int L$, these give
$\chi,\mathcal V=O(1)$.  Smoothness of the physical chart then gives
\begin{equation}
 |b|\le C\delta,\quad |a_T|\le C,\qquad
 |p-\pi\delta|+|h(\phi)-h_{\rm in}|\le Ce\delta s.
 \label{con:coefficient-drift}
\end{equation}
In particular $p/\delta$ is allowed to vary by $O(e)$.

First stop at $U=u_a$, positive-time force zero, a normal or mechanical
box boundary, or the reference section bracket endpoint.  On this
interval $U\ge u_a$, and (F4) gives
$\delta^{-1}F_\theta(U/\tau)\le
C\delta^{-1}e^{-\alpha u_a/\tau}\le C\tau^2/\delta=Ce^2\le Ce$.
Subtract \eqref{con:reference} from \eqref{con:actual-compression}.
The error has zero initial data and satisfies
\begin{equation}
 \begin{gathered}
 E_Y'=E_{V_c},\quad \sigma E_Z'=E_Y-2E_Z,\quad
 E_{V_c}'=\pi E_L+R_V,\quad E_U'=-h_{\rm in}E_L+R_U,\\
 R_V=\frac{p-\pi\delta}{\delta}L+\frac e\delta b
       +\frac{\mu B}{\delta}F_\theta(U/\tau)L,\qquad
 R_U=ea_T-(h(\phi)-h_{\rm in})L+\mu AF_\theta(U/\tau)L.
 \end{gathered}
 \label{con:compression-error}
\end{equation}
Both residuals are $O(e)$.  The zero-data filter has mass at most
$1/2$, so integration and Gronwall give
\begin{equation}
 |E_{V_c}(s)|+|E_U(s)|\le Ces,\qquad
 |E_Y(s)|+|E_Z(s)|+|E_L(s)|\le Ces^2.
 \label{con:initial-error}
\end{equation}
The quadratic force error and \eqref{con:initial-positive} exclude
an immediate force zero; the later reference force margin excludes
the others.  The estimates strictly improve the normal boxes,
and \eqref{con:mechanical-increments} improves the mechanical boxes.
The reference section bracket and its negative slip slope force the
first descending section $U=u_a$.  There is no earlier $M_{\rm sl}$ section,
since $u_a/\tau>M_{\rm sl}$ for small $\delta$.

Let $s_a,\bar s_a$ be these actual and reference section times.
They satisfy
\begin{equation}
 |s_a-\bar s_a|\le Ce,\qquad
 E_a=(Y,Z,V_c,U)(s_a)-
       (\bar Y,\bar Z,\bar V_c,\bar U)(\bar s_a)=O(e),
 \qquad E_{a,U}=0.
 \label{con:section-error}
\end{equation}
Only the reference is moved in this comparison.  Its derivative is
bounded: for example
$\bar Z'=-\zeta e^{-2s/\sigma}
+\sigma^{-1}\int_0^s e^{-2(s-t)/\sigma}\bar V_c(t)\,dt=O(1)$.

\paragraph{Step 3: The descending friction tail.}
On the tail compare times $s_a+\rho$ and $\bar s_a+\rho$.
Put $w=U/\tau$ and $G_s=H_\theta(w,X)L-ea_T$.
Stop at the first $w=M_{\rm sl}$, $L=\ell_a$, $G_s=g_a$, a box boundary,
or the reference endpoint $\bar U=-u_a/2$.  Choose
$0<g_a<\kappa_H\ell_a/4$ and then reduce $\delta$.  On every stopped
prefix,
\begin{equation}
 U'=-G_s,\quad w'=-G_s/\tau,\qquad
 G_s\ge(\kappa_H/2)L,\quad G_s\ge g_a,
 \quad L/G_s\le2/\kappa_H.
 \label{con:actual-tail}
\end{equation}
Integrating with respect to $U$ first gives $\int L\,d\rho\le C$.
Consequently \eqref{con:mechanical-increments} and
\eqref{con:coefficient-drift} hold with constants independent of the
as yet unchosen unweighted $V_c$ box.

For each stopped endpoint $\rho_b$, write $w_b'=w(\rho_b)\ge M_{\rm sl}$.
The actual decreasing slip gives
$d\rho=-\tau\,dw/G_s$ and
$\rho_b-\rho(w)\le(\tau/g_a)(w-w_b')$.  Hence (F4) gives the two
actual tail moments
\begin{equation}
 \begin{aligned}
 I_0&=\int_0^{\rho_b}F_\theta(w)L\,d\rho
       \le C\tau\int_{w_b'}^\infty e^{-\alpha w}\,dw\le C\tau,\\
 I_1&=\int_0^{\rho_b}(\rho_b-\rho)F_\theta(w)L\,d\rho
       \le C\tau^2\int_{w_b'}^\infty(w-w_b')e^{-\alpha w}\,dw
       \le C\tau^2.
 \end{aligned}
 \label{con:tail-moments}
\end{equation}
These are changes of variable on the actual orbit.

Let $S_\sigma$ be the propagator of the homogeneous system in
\eqref{con:compression-error}, ordered as $(Y,Z,V_c,U)$.
The filter mass and an integral Gronwall estimate give
$\|S_\sigma(\rho)\|\le C$ on the fixed horizon.  For
the column $v^{(c)}=S_\sigma(\rho)e_{V_c}$ the same argument, now with
$v_Y^{(c)}(0)=v_Z^{(c)}(0)=v_U^{(c)}(0)=0$, gives
$|v_{V_c}^{(c)}|\le C$ and
$|v_Y^{(c)}|+|v_Z^{(c)}|+|v_U^{(c)}|\le C\rho$;
also $S_\sigma(\rho)e_U=e_U$.  Duhamel's formula retains
$S_\sigma(\rho_b)E_a$.  Applying these bounds to the two residual
rows in \eqref{con:compression-error} gives
\begin{equation}
 \begin{aligned}
 |E_{V_c}|&\le C(e+\delta^{-1}I_0)\le C(e+e/\sigma),\\
 |\sigma E_{V_c}|+|E_Y|+|E_Z|+|E_U|
 &\le C(e+\sigma\delta^{-1}I_0+\delta^{-1}I_1+I_0)\le Ce.
 \end{aligned}
 \label{con:weighted-tail}
\end{equation}
Here $\sigma\tau/\delta=e$, $\tau^2/\delta=e^2$, and $\tau\le e$.
The first line need not be $O(e)$.

The second line strictly improves the force and normal boxes, and
\eqref{con:actual-tail} improves the descent margin.  The preceding
impulse bound improves the mechanical boxes.  We may now choose the
unweighted velocity radius larger than the reference bound plus
$C(1+C_0)$, and reduce $\delta$ last.  This excludes its boundary
by the first line.  If no $M_{\rm sl}$ event occurred by the reference endpoint,
actual $U\ge\tau M_{\rm sl}>0$ would contradict $\bar U=-u_a/2$ and
$E_U=O(e)$.  Thus $M_{\rm sl}$ is the first descending crossing and no release
has preceded it.  The reference negative slip slope gives
$|s_{\rm sl}-T_\sigma|\le C(e+\tau)$; moving only that reference gives
\eqref{con:M-admission}.  All preceding estimates hold on every
prefix, and $\int\lambda_N\,dt=\delta\int L\,ds\le C\delta$.

\paragraph{Step 4: Projection through the first release.}

By \eqref{con:inherited-amplitude} and Lemma~\ref{con:normal-opening},
$N_0-A_sn_\#$ and its derivative are $O(\sigma)$ on a common finite
interval.  Choose $r_-<\rho_\#<r_+<R_*$ so that, after reducing
$\delta$, the inherited reference satisfies
\begin{equation}
 \begin{gathered}
 L_0\ge m_{\rm pre}>0\ (0\le\rho\le r_-),\quad
 L_0(r_+)\le-m_{\rm end}<0,\quad
 L_0'\le-m_{\rm der}<0\ (r_-\le\rho\le r_+),\\
 Y_0(R_0)=Z_0(R_0)\le-m_Y,
 \quad V_0(R_0)\ge m_V,\quad
 V_0(R_0)+Z_0(R_0)\ge m_{\rm der}.
 \end{gathered}
 \label{con:reference-release-margins}
\end{equation}
This proves that $R_0$ is the first root and $R_0=\rho_\#+O(\sigma)$.

In actual release variables the exact positive-contact equations are
\begin{equation}
 \begin{aligned}
 Y'&=\widehat V,& Z'&=Y-2Z,\\
 \widehat V'&=eb+(C-\mu B\Phi_\theta(w))L,&
 w'&=a_T+e^{-1}(B-\mu A\Phi_\theta(w))L,
 \qquad w=u_r/\eps.
 \end{aligned}
 \label{con:actual-release}
\end{equation}
At $\rho=0$, $\widehat V=\sigma V_c(s_{\rm sl})$ and $w=M_{\rm sl}$.
Stop at first force zero, $R_*$, or a normal or mechanical boundary.
For $w\ge M_{\rm sl}$ the friction term is nonpositive, so
$w\le W_*:=M_{\rm sl}+R_*\sup|a_T|$.  Since $B,a_T>0$ and
$\Phi_\theta(0)=0$, a common $w_b\in(0,M_{\rm sl})$ satisfies
$B-\mu A\Phi_\theta(w)\ge B_{\min}/2$ on $[0,w_b]$.
The vector field there points upward, and hence $w_b\le w\le W_*$
throughout the stopped interval.

Set $\vartheta=B/A$ and $P=\widehat V-e\vartheta w$.
The variable coefficients must be differentiated before estimating.
The identity $C-B^2/A=1/2$ gives exactly
\begin{equation}
 \widetilde N=(Y,P,Z),\qquad
 \widetilde N'=A_D\widetilde N+eF,\qquad
 F=(\vartheta w,b-\vartheta a_T-\vartheta'w,0),
 \quad \vartheta'=\eps\vartheta_\phi\dot\phi.
 \label{con:projection}
\end{equation}
No derivative of $w$ remains.  On the stopped boxes $F$ is bounded.
The original mechanical equation also gives, for every such prefix,
\begin{equation}
 |q-q_{\rm sl}|\le C\eps,\qquad
 |\dot q-\dot q_{\rm sl}|\le C\eps+
 C\delta^{3/2}\int_0^\rho L\,dt
 \le C(\eps+\delta^{3/2}),
 \label{con:release-mechanics}
\end{equation}
which strictly improves the mechanical box.
The initial error is
$\widetilde N(0)-N_0(0)=
\{(Y,\widehat V,Z)(0)-N_0(0)\}-(0,e\vartheta(0)M_{\rm sl},0)=O(e)$.
Duhamel's formula in \eqref{con:projection}, followed by
$\widehat V=P+e\vartheta w$, proves the first line of
\eqref{con:release-comparison} on every stopped prefix.  For the force
derivative use $L'=Y-2Z-P-e\vartheta w$, which again contains no $w'$.
This strictly improves the normal box.

The first margin in \eqref{con:reference-release-margins} excludes
an earlier force zero before $r_-$.  In the stopped bracket
$L'\le-m_{\rm der}/2$; if positive contact persisted to $r_+$, its
negative reference endpoint would give a contradiction.  The actual
first zero is therefore in $(r_-,r_+)$, and
$|L_0(R_\eps)|\le Ce$ gives $|R_\eps-R_0|\le Ce$.
Moving only the reference proves the remaining event comparison.
At this zero $Y=Z$ and
$-L'=\widehat V+Z$, so
$\dot g=v+z/\eps=\delta^{3/2}(\widehat V+Z)>0$.
This proves \eqref{con:opening}, including the negative memory term.

\paragraph{Step 5: Physical reconstruction and leading increments.}
The physical reconstruction in \eqref{con:physical-reconstruction}
uses zero physical displacement and memory, while
\eqref{con:inherited-release} retains their complete scaled values.
We have $u_{r,L}=\eps w=O(e\delta^{3/2})$ and
$y_L=z_L=O(\eps\delta^{3/2})=O(e\delta^3)$.
The $p$ drift through release is $O(\eps)$, and the common inverse
chart gives
$|\Delta b_L|\le C(|\Delta p_L|+|\Delta u_0|+|\Delta y_L|+
|\Delta v_L|)$, with $\Delta u_0=\Delta u_r$.
Thus \eqref{con:physical-errors} follows, with $b_L$ evaluated at
the outgoing state in both passages.  The independent section and
release clocks give \eqref{con:release-clock}.

The compression increments, its bounded horizon and impulse, and
\eqref{con:release-mechanics} prove every estimate in
\eqref{con:prefix}; the release impulse is $O(\delta^{3/2})$.
The velocity bound uses $|v|\le C\delta^2$ in compression and the
release scale thereafter.  The full state includes the continuous
memory coordinate at the joining point.  Finally
$\dot p=p_\phi\omega=-a+O(\delta)$ throughout this same prefix.

Integrating the compression slip row, using $U(s_{\rm sl})=\tau M_{\rm sl}$,
\eqref{con:tail-moments} and \eqref{con:coefficient-drift}, gives
\begin{equation}
 \int_0^{s_{\rm sl}}L\,ds=S/h_{\rm in}+O(e),\qquad
 \int_{t_{\rm in}}^{t_{\rm sl}}\lambda_N\,dt
   =\delta S/h+O\{\delta(e+\delta)\}.
 \label{con:compression-impulse}
\end{equation}
The discrepancy from saturated impulse is $\delta I_0=O(\delta\tau)$.
The saturated direction at $G$ is $(-5/4,1)^T$ in generalized
velocity, and $Q_{\rm sat}b(G)=a\beta$.
The free-force impulse is $O(e\delta)$ and the release impulse is
$O(\delta^{3/2})$.  These facts prove the velocity and $b$ rows of
\eqref{con:leading-data}.  Integrating the configuration and $p$ rows,
with $s_{\rm sl}=\xi_s+O(\sigma)$ and release duration $O(\eps)$, proves
their rows.  The normal-state formulas in \eqref{con:leading-data} follow from
\eqref{con:inherited-amplitude} and \eqref{con:release-comparison}.
All constants, boxes and event brackets are chosen after the fixed
classes and $C_0$, and the common scale cutoff is reduced last.
\end{proof}

\subsection{Parameter dependence of the contact reference}

To differentiate the limiting return, we let the initial displacement
and memory of the linear contact reference vary independently.

\begin{lemma}[Parameter dependence of the limiting contact]
\label{con:reference-regularity}
In \eqref{con:reference}, replace $h_{\rm in}$ by an independent
$h_0$ near $h$ and the initial data by $(y_0,z_0,-Q,S)$.
On a neighborhood of the compact data
$(0,\pi,h,0,0,Q,S)$, the first arrest, its inherited normal state,
and the first inherited release have $C^2$ extensions in
$(\sigma,\pi,h_0,y_0,z_0,Q,S)$ to $\sigma=0$, with uniform bounds
through order two and continuous dependence on $\lambda$.
If $\mathcal Q=V_0(R_0)$ denotes the release velocity, then
\begin{equation}
 \mathcal Q=
 \nu_\#\sqrt{QS/h_0-\pi S^2/(2h_0^2)}
       +O(\sigma+|y_0|+|z_0|).
 \label{con:reference-amplitude-regularity}
\end{equation}
The same error bound holds after one derivative in
$(\pi,h_0,Q,S)$; derivatives in $(\sigma,y_0,z_0)$ are uniformly
bounded.  Compression on this auxiliary neighborhood is the signed
linear continuation: its force need not be positive near time zero.
Its arrest is the first zero of $U$, and the subsequent inherited
release is the first zero of its normal force.  No physical closing
condition is imposed on the auxiliary initial data.
\end{lemma}
\begin{proof}
Set $R=\bar Z-\bar Y/2$.  The normal compression rows become
\[
 (Y,V,R)'=(V,-\pi Y/2+\pi R,-V/2-2R/\sigma).
\]
A slow invariant plane $R=mY+nV$ is determined by
\[
 2m+\sigma\pi n(m-1/2)=0,\qquad
 2n+\sigma(m+\pi n^2+1/2)=0.
\]
The implicit function theorem gives analytic $m,n$, uniformly for
$\pi$ in a neighborhood of its compact range, with
$m=O(\sigma^2)$ and $n=-\sigma/4+O(\sigma^3)$.
On this plane the matrix is
$B_s=\left(\begin{smallmatrix}0&1\\-\pi/2+\pi m&\pi n\end{smallmatrix}\right)$.
The remaining root is $\nu/\sigma$, where
$\nu=-2+\sigma^2g$ and
$\nu^2g+\pi(\nu+1)=0$.  Its eigenvector, normalized in the $R$ row,
is $(-2\sigma^2g/\nu,-2\sigma g,1)^T$.
Together with $(1,0,m)^T,(0,1,n)^T$ these columns form an analytic
invertible matrix at $\sigma=0$.
The exact solution thus consists of the analytic slow block
$e^{B_ss}$ and a fast term proportional to $e^{\nu s/\sigma}$,
with $\nu<-1$.  On every late interval $s\ge s_*>0$, that fast term
and its derivatives of any fixed order extend flatly to zero.
The slip row is integrated without dividing by $\pi$:
$U'=h_0Y/2-h_0R$.  Its fast contribution contains
$(\sigma/\nu)(e^{\nu s/\sigma}-1)$, which has the same smooth late
extension.  No diagonalization of the slow block is required at
$\pi=0$.

At zero scale and $y_0=z_0=0$, the arrest has a common positive time
and negative slip slope, by \eqref{con:amplitude}.  Its slip stays
positive before a fixed arrest bracket.  Uniform convergence of the
integrated solution on the whole preceding interval preserves this
first-event property for the small independent perturbations.
The implicit function theorem on the late bracket gives a $C^2$
arrest and inherited state.  Lemma~\ref{con:normal-opening} and the
linear formula $e^{A_D\rho}N_0(0)$ give the same conclusion for its
first release, with the positive-force segment and transverse bracket
preserved.  At $\sigma=y_0=z_0=0$ its velocity is exactly the first
term of \eqref{con:reference-amplitude-regularity}.
The bounded second derivatives and the mean value formula give both
the value and the stated first-derivative bounds.
\end{proof}

For subsequent prepared touchdowns set
\begin{equation}
 \pi=-a,\quad Q=r_0W,\quad S=d_\lambda U^r,\quad
 r_0=c_\lambda/3,\qquad
 W=1+\delta^{1/3}\eta_W,\quad
 U^r=1+\delta^{1/3}\eta_U.
 \label{con:native-contact}
\end{equation}
The prepared cone retains its complete dependent acceleration
coordinate, $|\zeta|\le Z_*$ and $r_0W-\zeta\ge\gamma_0$;
its inverse is taken at $u_0=d_\lambda\delta U^r+r$ with
$0\le r\le A_{\rm off}\delta^{4/3}$.
These hypotheses give \eqref{con:input}.  The reference is exactly
\eqref{con:reference} with its finite $h_{\rm in}$ and complete
initial memory, so \eqref{con:physical-errors} applies to this same
input and offset.
For the tuned limiting reference define
\begin{equation}
 q_\lambda=\nu_\#
 \sqrt{r_0d_\lambda/h+ad_\lambda^2/(2h^2)}>0.
 \label{con:q-lambda}
\end{equation}
The analytic inverse of $p(\phi_\delta)=-a\delta$ gives
$h_\delta=h+O(\delta)$ and $\delta h_\delta'=O(\delta)$.
Lemma~\ref{con:reference-regularity}, restricted to
$\sigma=\delta^{1/2}$, $h_0=h_\delta$,
$y_0=z_0=\sigma\zeta$ and \eqref{con:native-contact}, gives
\begin{equation}
 \mathcal Q=q_\lambda+O_{C^1_*}(\delta^{1/3}),\qquad
 |D_\xi\mathcal Q|\le C\delta^{1/3},\qquad
 \xi=(\zeta,\eta_U,\eta_W).
 \label{con:weighted-contact}
\end{equation}
Here $C^1_*$ uses $\delta\partial_\delta$ and $D_\xi$ on any fixed
bounded auxiliary fibre set.  The chain includes both initial memory
coordinates: its radial part is
$(\sigma/2)\partial_\sigma+
\delta h_\delta'\partial_{h_0}+
(y_0/2)\partial_{y_0}+(z_0/2)\partial_{z_0}$, together with
$\delta\partial_\delta Q$ and $\delta\partial_\delta S$ from
\eqref{con:native-contact}.  Thus the reference is defined on the
open neighbourhood needed for fibre differentiation, even where
physical closing is not imposed.  The free-flight construction below
will establish the complete weighted return factorization.

The first physical touchdown uses the wider offset range
$\delta=\sqrt\eps$, $r=k\delta$. Its complete input and the
strict signed arrest margin are verified in
Proposition~\ref{prep:first-flight}; later prepared inputs are
verified in Section~\ref{sec:detuned-return}.
